\scanpage{3}

\deleted{\unsure{Notamment|Néanmoins} on voudrait que}

\deleted{\[
\underline{\mathrm{Ar}}(C)\xrightarrow{p}C\times C
\]}

\deleted{ait des propriétés du type ``fibration'', \deleted{\unsure{ou|on}} Or

1) au sens catégorique, ce n'est ni une fibration, ni une cofibration -- et $=$ les \deleted{\unsure{projets|projections|propriétés}} composés $\underline{s}\circ p,\underline{b}\circ p\colon\underline{\mathrm{Ar}}(C)\longrightarrow C$ ne le sont pas.}

\deleted{2) De $=$, \deleted{\unsure{ni}} $p$ \deleted{\unsure{n'est}} n'est ni cohom. propre, ni coh. lisse. Peut-être au moins $\deleted{\unsure{s}}\ \underline{s}\circ p$ et $\underline{s}\circ b$ \deleted{le} sont-ils \deleted{\unsure{le}} l'un ou l'autre ? Mais ce sont plutôt les propr. de $p$ lui-$=$ qui m'importent !}

\deleted{3) Il semblerait que $p$ soit une $W_\infty$-fibration. (``fibration cohomologique'') C'est une chose à vérifier.}

\deleted{\uline{Question} Les $W_\infty$-fibrations, et la notion correspondante \inserted{(par le L.L.P)} de cofibration, \deleted{\unsure{font}} font-ils une \deleted{des} structure de Quillen sur $(Cat)$ ?}

% archive sheet number: 2
